% Reproducible source for the BIBM 2026 Doctoral Forum cover letter.
% Mirrors manuscript/cover_note.md. Compile: pdflatex cover_bibm2026.tex
\documentclass[11pt]{article}
\usepackage[margin=1in]{geometry}
\usepackage{parskip}
\pagenumbering{gobble}
\begin{document}

\begin{center}
{\large\textbf{IEEE BIBM 2026 --- Doctoral Forum Cover Note}}
\end{center}

\vspace{1em}

Dear Dr.\ Yi He and Dr.\ Xingquan Zhu,

I am pleased to submit ``Spiking Reservoir State-Space Encoding for Perturbation-Characterized Affective EEG Signal Analysis'' for the IEEE BIBM 2026 Doctoral Forum. The paper evaluates a fixed leaky integrate-and-fire (LIF) reservoir with a six-bin spike-count (BSC6) code as a nonlinear temporal encoding operator for affective EEG/ERP observations. Using subject-grouped validation against classical band-power and ERP-window baselines, it reports balanced accuracy, macro-F1, and macro one-vs-rest AUC, and characterizes the encoding under amplitude noise, temporal jitter, and channel dropout, separating representation-level corruption from full raw-signal recomputation.

The result is deliberately bounded. The reservoir embedding is not presented as a uniformly stronger static classifier than classical features; it is comparable on clean separability but exhibits a distinct perturbation-response profile, retaining balanced accuracy more strongly under representation-level amplitude noise and channel dropout while remaining sensitive to temporal misalignment. Controls --- a dimensionality-matched random projection, a label-permutation null, and a single-bin spike-count ablation --- indicate that this robustness reflects the fixed reservoir state-space expansion rather than feature dimensionality or temporal bin granularity. Reporting what a spiking representation preserves or suppresses under controlled corruption, rather than a single clean-accuracy number, gives a reproducible operating-regime account for more auditable EEG/ERP pipelines.

The work fits the Doctoral Forum: it reframes a spiking reservoir as a perturbation-characterized temporal measurement operator for affective EEG/ERP rather than as another endpoint classifier, and its evaluation --- subject-grouped validation, classical baselines, and explicit confound controls --- is designed for soundness and clarity. The first author is currently enrolled in the Ph.D.\ program in Electrical and Computer Engineering at Stony Brook University, is the primary contributor, and will present the work if accepted; a one-page CV of the first author is appended as required by the call. This submission is intentionally focused on the reservoir temporal-encoding component of the broader doctoral work.

\vspace{1em}

Sincerely,

Andrew Lane

\end{document}
